% \begin{savequote}
% {\tiny Ein Mann ging von Jerusalem nach Jericho hinab und wurde von
% Räubern überfallen. Sie plünderten ihn aus und schlugen ihn nieder; dann gingen sie weg und ließen ihn 
% halbtot liegen. Zufällig kam ein Priester denselben Weg
% herab; er sah ihn und ging weiter. Auch ein Levit kam zu
% der  Stelle; er sah ihn und ging weiter. Dann kam ein Mann
% aus  Samarien, der auf der Reise war. Als er ihn sah, hatte
% er  Mitleid, ging zu ihm hin, goss Öl und Wein auf seine
% Wunden  und verband sie. 
% 
% Dann hob er ihn auf sein Reittier,
% brachte  ihn zu einer Herberge und sorgte für ihn. Am
% anderen Morgen  holte er zwei Denare hervor, gab sie dem
% Wirt und sagte:  Sorge für ihn, und wenn du mehr für ihn
% brauchst, werde  ich es dir bezahlen, wenn ich wiederkomme.
% Was meinst du:  Wer von diesen dreien hat sich als der
% Nächste dessen erwiesen,  der von den Räubern überfallen
% wurde?  
% 
% Der Gesetzeslehrer antwortete: Der, der
% barmherzig\index{Barmherzigkeit} an  ihm gehandelt hat. Da
% sagte Jesus zu ihm: Dann geh und handle genauso!
% 
% Lk 10,25-37}
% % \qauthor{{\tiny Lk 10,25-37}}
% \end{savequote}

% \clearpage
\def\ChapterCode{ASB}
\chapter
[Absaugung, Sauerstoff, Beatmung]
{Absaugung, Sauerstoff, Beatmung}
\ChapterIntro
{%
~~~
}
{%
\begin{description}
  \item[Maintainer:] Sebastian Gabriel
  \item[Autoren:] Diverse
  \item[Reviewer:] Standard-Reviewprozess
  \item[Version:] {\DocVersion} ({\DocVersionInfo})
  \item[SHA1:] {\DocGitCommit}
\end{description}

{\TxTeilkapitelAASS}
}

\clearpage
\section{Absaugung}

\label{topic:absaugkatheter}
\label{topic:absaugung-taetigkeit}
\index{Absaugung}
\index{Absauggeräte}
\index{Absauggeräte!Accuvac{\texttrademark}}
\index{Absauggeräte!Laerdal Suction Unit{\texttrademark}}
\index{Accuvac{\texttrademark}}
\index{Laerdal Suction Unit{\texttrademark}}
\index{Absaugkatheter}
\label{topic:absaugung}
\label{topic:absauggeraete}

% \Text{{\TxQuerverweise}}{%
% \begin{inparaitem}
%   \item Absaugkatheter: \prettyref{topic:absaugkatheter};
%   \item Absaugtechnik: \prettyref{topic:absaugung-taetigkeit}
% \end{inparaitem}
% }{}{}{}{}

\Text{Einleitung}
{Blut oder Schleim kann vom Fachpersonal mittels
Absauggeräten aus dem Mund-Rachen-Raum abgesaugt werden. Es
gibt unterschiedliche Modelle; grundsätzlich wird nach der
Betriebsart unterschieden zwischen \EE{elektrischer
Absaugeinheit}, \EE{manueller Absaugpumpe} und als
Sonderfall dem \EE{Oro-Sauger} oder Schleimabsauger.

Ein \EE{Absaugkatheter} ist jenes Zubehör zum Absauggerät,
mit welchem die Sekrete aus Mund oder Nase abgesaugt werden
können. Es ist ein flexibler, ca. 50\CM* langer Schlauch,
der direkt an die Absaugpumpe angeschlossen wird (siehe
Abbildung in \prettyref{tab:absaugvorrichtungen}). Die
meisten Absaugkatheter haben auf Patientenseite eine kleine
Öffnung, die das Festsaugen an dem Patienten verhindert. Die
Absaugkatheter sind steril verpackt und müssen auch
dementsprechend behandelt werden. Die farbige Kennung gibt
den Durchmesser des Absaugkatheters an. \Z{Die Farbcodierung
kann \prettyref{tab:absaugkatheter} entnommen werden.}}
{}
{}
{}
{}



\Layout{57}{}
{}{}
{\Z{%
\begin{tabulary}{\linewidth}{LCCCCCCCCC}\toprule\addlinespace
\noindent\TabHeading{Farbe}
&
\noindent\TabHeading{grau}
&
\noindent\TabHeading{hellgrün}
&
\noindent\TabHeading{blau}
&
\noindent\TabHeading{schwarz}
&
\noindent\TabHeading{weiß}
&
\noindent\TabHeading{grün}
&
\noindent\TabHeading{orange}
&
\noindent\TabHeading{rot}
&
\noindent\TabHeading{gelb}
\\\addlinespace\midrule
Durchmesser in mm
& 1,7 & 2,0 & 2,7 & 3,3 & 4,1 & 4,7 & 5,3 & 6,1 & 6,6
\\\addlinespace
Größe in Charrière 
& CH\,5 & CH\,6 & CH\,8 & CH\,10 & CH\,12 & CH\,14 & CH\,16
& CH\,18 & CH\,20 
\\\addlinespace\bottomrule
\end{tabulary}
}
}
{Farbcodierung: Absaugkatheter\label{tab:absaugkatheter}}{}


\Text{Oro-Sauger}{
Der Oro-Sauger ist ein simples Gerät, welches zum Absaugen
von Schleim von \E{Neugeborenen} eingesetzt wird. Er
besteht aus einem Absaugkatheter, welcher mit einem Behälter
verbunden ist. Von dem Behälter geht ein weiterer Schlauch
ab, welcher vom Helfer in den Mund genommen wird. Der
Oro-Sauger wird \E{nicht} an ein Absauggerät angeschlossen,
der Helfer muss selbsttätig mit dem Mund
(\Lang{lat.} \I{oro}) saugen. Dadurch gelangt das
Sekret aus dem Absaugkatheter in den Behälter. 
Der Katheter und der andere Schlauch müssen dabei 
\E{von oben} den Behälter verlassen,
um zu verhindern, dass Sekret in den Mund des Helfers
gelangt.
Bei den meisten Fabrikaten kann der Behälter abgeschraubt
und verschlossen werden.

}{}{}{}{}





% \opt{ORGASBFLO}{\Customized{%
% \begin{description}
%   \item[\CustAsbFlo:] Es werden folgende Geräte
% eingesetzt:
% 
% \begin{tabulary}{\linewidth}{LLLL}\toprule\addlinespace
% \noindent\TabHeading{Hersteller}
% &
% \noindent\TabHeading{Produktfamilie}
% &
% \noindent\TabHeading{Typ}
% &
% \noindent\TabHeading{Bemerkungen}
% \\\addlinespace\midrule
% Weinmann
% &
% \noindent\T{Accuvac{\texttrademark}}
% &
% Rescue
% &
% rote Deckplatte
% \\\addlinespace
% &
% &
% Basic
% &
% grüne Deckplatte
% \\\addlinespace
% Laerdal
% & 
% Laerdal Suction Unit{\texttrademark}
% &
% &
% auslaufend
% \\\addlinespace
% div. Hersteller
% &
% Handabsaugpumpe
% &
% RES-Q-VAC
% &
% \\\addlinespace
% div. Hersteller
% &
% Orosauger
% &
% &
% \\\addlinespace\bottomrule
% \end{tabulary}
% \end{description}
% }}%Ende opt

\Layout{57}{}{}{}
{
\begin{tabulary}{\linewidth}{LLL}\toprule\addlinespace
\noindent\TabHeading{Elektrische Absaugeinheit}

Älteres Modell von der Fa. Laerdal (Laerdal Suction Unit{\texttrademark})

Neueres Modell von der Fa. Weinmann (Accuvac{\texttrademark})
&
\vspace{0.1pt}
\includegraphics[width=\linewidth]{./images/pallinger-christoph-aass/LSU_32831_export_00800.jpg}
&
\vspace{0.1pt}
\includegraphics[width=\linewidth]{./images/pallinger-christoph-aass/Accuvac_32820-AASS-0112mm.jpg}
\\\midrule
\noindent\TabHeading{Manuelle Absaugpumpen}

Handabsaugpumpe

Fußabsaugpumpe
&
\vspace{0.1pt}
\includegraphics[width=\linewidth]{./images/geraete/handabsaugung.jpg}
&
\vspace{0.1pt}
\includegraphics[width=\linewidth]{./images/geraete/fussabsaugung1.jpg}
\\\midrule
\noindent\TabHeading{Oro-Sauger}

Schleimabsauger für Neugeborene
&
\vspace{0.1pt}
\includegraphics[width=\linewidth]{./images/geraete/orosauger.jpg}
\\\midrule
\noindent\TabHeading{Absaugkatheter}

zum Einführen in Mund oder Nase
&
\vspace{0.1pt}
 \includegraphics[width=4cm]{./images/geraete/absaugkatheter.jpg} \\\bottomrule
\end{tabulary}
}
{Verschiedene Absaugvorrichtungen.\label{tab:absaugvorrichtungen}}
{}

{\TakeNotesBoxLarge}

% \begin{table}[H]
% \begin{WidePar}
% \Captionof{table}{}
% 
% 
% \end{WidePar}
% \end{table}

\TextSlide{Funktionskontrolle}
{%
Jede Funktionskontrolle umfasst die Überprüfung der
\EE{Vollständigkeit} des Materials, \EE{Saugfunktion} bei
maximaler und minimaler Leistung und \EE{Dichtheit}. Sie muss
zumindest bei Dienstbeginn durchgeführt werden.
Bei Geräten, welche über Akkus verfügen, muss die
\EE{Akku-Ladung} und Saugfunktionsprüfung \E{ohne Anschluss
an eine externe Stromquelle} erfolgen, d.h. das Gerät muss
für die Funktionsprüfung von einer Ladehalterung abgenommen
werden. Die Ladestandsanzeige und evtl. vorhandene
\EE{Warnleuchten} müssen kontrolliert werden.
Die vom jeweiligen Hersteller vorgegebenen,
typenspezifischen Anweisungen müssen beachtet werden. 
}{
\begin{itemize}
   \item Prüfung von
   \begin{itemize}
      \item Vollständigkeit
      \item Saugleistung
      \item Dichtheit
      \item Ggfs. Akku-Ladung
      \begin{itemize}
         \item Ohne externe Stromquelle!
      \end{itemize}
      \item Warnleuchten
   \end{itemize}
   \item Herstellerhinweise beachten!
\end{itemize}
}{}{}{}



\Layout{37}{Weinmann AccuVac{\protect\texttrademark}}
{}
{}
{./images/pallinger-christoph-aass/Accuvac_32820-AASS-0176mm.jpg}
{Weinmann AccuVac{\protect\texttrademark}}
{Ch. Pallinger}




\TextSlide{Absaugbereitschaft}
{%
\index{Absaugbereitschaft}
Wenn \EE{Absaugbereitschaft} hergestellt werden muss, wird
aus hygienischen Gründen nur der farbige Trichter des
Katheters ausgepackt und an den Konnektor bzw. den
\T{Fingertip} der Absaugpumpe angeschlossen. Der restliche
Teil des Katheters bleibt so lange in der Verpackung bis mit
der Absaugung begonnen
wird.\bigskip}{
\begin{itemize}
  \item Verpackung nur am farbigen Ende des Katheters öffnen
  \item An Absaugpumpe anschließen
  \item Katheter bleibt in Hülle
\end{itemize}
}
{}
{}
{}

\TextSlide{Absaugtechnik}{%
Damit der
Absaugkatheter nicht zu tief eingeführt wird, muss die
einzuführende Katheterlänge durch Abmessen der Distanz
zwischen \EE{Mundwinkel und Ohrläppchen} bzw. zwischen \EE{Nasenspitze
und Ohrläppchen} bestimmt werden. Es darf \EE{maximal auf Sicht} abgesaugt
werden!

Der Katheter wird \EE{ohne Sog eingeführt}, dazu wird entweder der Katheter
abgeknickt  oder der Fingertip offen gelassen. Unter kreisenden Bewegungen wird
dann auf Sicht abgesaugt und der Katheter herausgezogen. Wenn sich der Katheter
festsaugt, muss der Fingertip geöffnet oder der Katheter abgesteckt werden.
Der Vorgang wird nach Bedarf wiederholt.

\Cave{Es darf nur auf Sicht abgesaugt werden!}

Bei Verdacht auf ein \EE{Schädel-Hirn-Trauma}
(\FullPrettyRef{topic:sht}) darf \EE{nicht über die Nase}
abgesaugt werden. 
}{
\begin{itemize}
  \item Tiefe:
  \begin{itemize}
    \item Max. auf Sicht.
    \item Mund: Mundwinkel -- Ohrläppchen.
    \item Nase: Nasenspitze -- Ohrläppchen.
  \end{itemize}
  \item Ohne Sog einführen
  \begin{itemize}
    \item Katheter abknicken, Fingertip offen
  \end{itemize}
  \item Unter kreisenden Bewegungen saugen und herausziehen
\end{itemize}
}{}{}{}

\begin{table}[H]
\TableBaseModification
\Captionof{table}{Bilderserie: Absaugtechnik}
\begin{tabularx}{\linewidth}{XX}
\includegraphics[width=\linewidth]{./images/pallinger-christoph-aass/Anton_32999_export.jpg}
\Captionof{figure}{Katheterlänge bei Absaugung des Mundes \SourcePicture{Ch. Pallinger}}
&
\includegraphics[width=\linewidth]{./images/pallinger-christoph-aass/Anton_32990_export.jpg}
\Captionof{figure}{Katheterlänge bei Absaugung durch die Nase \SourcePicture{Ch. Pallinger}}
\end{tabularx}
\end{table}





\TextSlide{\TxGefahren}{%
In der Rachenhinterwand verläuft ein Nerv, welcher Einfluss auf den Kreislauf
hat (Vagus-Nerv). Wenn dieser gereizt wird kann es zu gefährlichen
\EE{Kreislaufstörungen} (Bradykardie, Hypotension) kommen. Der Vagus-Nerv ist
Teil des \E{Parasympathikus} und des vegetativen Nervensystems
(\FullPrettyRef{topic:parasymapthikus}). 
Weiters kann es zu \EE{Schleimhautverletzungen} und
Blutungen kommen. Ebenso kann ein \EE{Brechreiz} ausgelöst werden,
wenn der Katheter zu tief eingeführt wird. Bei Verdacht auf
ein Schädel-Hirn-Trauma (\FullPrettyRef{topic:sht}) darf aufgrund möglicher
Verletzungen der Schädelbasis nicht über die Nase abgesaugt
werden!

 }{
\begin{itemize}
  \item Reizung des Vagus-Nerv
  \item Gefährliche Kreislaufstörungen
  \item Schleimhautverletzungen, Blutungen
  \item Brechreiz
\end{itemize}
}{}{}{}


\needspace{0.5\textheight}
\section{Sauerstoff -- \texorpdfstring{\ce{O2}}{}}
\label{topic:sauerstoff}
\label{topic:o2}
\index{Sauerstoff}
\index{O2@\ce{O2}}

\Definition{\T{Sauerstoff} ist ein
lebenswichtiges Gas welches zu 21\PC* in unserer Atemluft
vorkommt und außerdem ein essentielles
Notfall\enquote{medikament} ist.
Das chemische Symbol lautet \T{\ce{O2}}.}

% \begin{center}
% \includegraphics[width=4.522cm,height=5.628cm]{./images/rippenspreitzer-02-fett.png}
% \Captionof{figure}{NEIN! \SourcePicture{Rippenspreitzer.de}}
% 
% \end{center}

\subsection{Technische Grundlagen}
% TODO : Bild Druckminderer
\TODO{GABS}{2010-08-04}{Bild: Sauerstofffl., Druckminderer fehlt.}{}
%
Sauerstoff ist ein farb- und geruchsloses Gas und wird in
Druckflaschen gelagert.

\TextSlide{Kennzeichnung}
{Die \EE{Farbe} von Gasflaschen ist
normiert\ZFootnote{Normen zur Kennzeichnung von Gasflaschen:
EN~1089-3, in der jeweiligen nationalen Umsetzung (z.\,B.
DIN EN~1089-3, ÖNORM EN~1089-3 (Verbindlichkeit gem.
Versandbehälterverordnung~2002 (VBV~2002)), \ldots).

Die Farbcodierung wurde in der Zeit von
\VonBis{1998}{2006}\ZFootnote{Die Umstellung auf das neue
Farbsystem gemäß ÖNORM EN 1089-3 musste bei allen anderen
Gasen bis zum \E{30. Juni 2006} abgeschlossen sein (\S\,26
Abs.\,3 VBV 2002 \cite{AUT-LAW-VBV2002-1}).} umgestellt. In
der Übergangszeit wurden Flaschen mit der neuen Codierung
mittels Aufbringung des Buchstabens \enquote{\EE{\sffamily
N}} (\enquote{neue Kennzeichnung}) gekennzeichnet. Dies hat
\E{nichts} mit dem chemischen Symbol für Stickstoff (\ce{N})
zu tun!}.
Flaschen mit Sauerstoffbefüllung werden an der
Flaschenschulter mit einem \EE{weißem Farbring}
gekennzeichnet.
Gase für den medizinischen Gebrauch werden zusätzlich durch
einen weißen Flaschenkörper gekennzeichnet.
Gasflaschen mit \EE{medizinischem Sauerstoff} sind daher
\EE{vollständig weiß}\ZFootnote{Gasflaschen für
\E{technischen} Sauerstoff haben zwar eine weiße
Flaschenschulter, jedoch einen \E{blauen Flaschenkörper}.
Technischer Sauerstoff darf Patienten nicht verabreicht
werden.}.
\cite{AUT-LAW-VBV2002-1,Igv2010} 
Weiters ist auf Druckgasflaschen ein
\EE{Gefahrengutaufkleber} angebracht.

\Fazit{Gasflaschen mit \EE{\E{medizinischem} Sauerstoff}
sind
\EE{vollständig weiß}.}

\Commentary[Gasflaschen / Farbcodierung / Umstellung]
{\S\,26 Abs.\,2 VBV 2002:
Bei Sauerstoff, Distickstoffmonoxid (Lachgas) und
oxidierenden Gasen erfolgt die Umstellung vom Farbsystem
gemäß \S\,15 Abs. 1 Z 2 lit. a für die industrielle Anwendung
ab dem 1. Jänner 2002 und bei solchen für medizinische
Zwecke ab dem 1. Jänner 2004.
\ParBugfix
Abs. 3: Die Umstellung auf das neue Farbsystem gemäß ÖNORM
EN 1089-3 hatte bei Acetylen bis zum 31. Dezember 2001 stattzufinden,
bei allen anderen Gasen muss sie bis zum 30. Juni 2006
abgeschlossen sein.

\cite{AUT-LAW-VBV2002-1}
}
\Commentary[Gasflaschen / Farbcodierung / EN~1089-3]
{Nach: \cite{Igv2010}:

Die verbindliche Kennzeichnung des Gaseinhalts erfolgt auf
dem Gefahrgutaufkleber.
Die Farbkennzeichnung ist nur für die Flaschenschulter
festgelegt,
 außer bei medizinischen Gasen. In diesem Fall ist der
 zylindrische Teil weiß.
Die Farbe des zylindrischen Flaschenmantels ist in der Norm
(bis auf medizinischen Gase) nicht festgelegt.
}

}
{
\begin{itemize}
  \item Medizinischer Sauerstoff: Vollständig weiß
  \item Gefahrengutaufkleber
\end{itemize}
}
{}
{}
{}

\Layout{27}{}
{}% Text
{}% Slides
{./images/pallinger-christoph-aass/O2_Flasche_32747-crop2-AASS-0176mm.jpg}% Bildpfad oder Objeckt
{Eine Sauerstoffflasche mit angeschlossenem
Druckminderer/Berieselungseinheit}% Beschreibung
{Pallinger}% Credits


\TextSlide{Druckflaschen}
{Die Druckflaschen haben unterschiedliche Füllgrößen
(Volumina), gängige Flaschengrößen sind 0,5, 1, 1,5, 2, 5
und
10\Liter*. Der Sauerstoff wird unter Druck gelagert, wodurch
die Menge des zur Verfügung stehenden Sauerstoffs um ein Vielfaches
gesteigert wird. Der Druck einer gefüllten Flasche beträgt
normalerweise ca. 200\Bar*. 

Jede Druckflasche verfügt über ein \T{Hauptventil} und
ein Standardgewinde. An diesem Gewinde ist entweder ein
weiterführender Druckschlauch oder der \E{Druckminderer}
einer \E{Berieselungseinheit} angeschlossen.

\bigskip }
{
\begin{itemize}
  \item Unterschiedliche Füllgrößen
  \item Druck ca. 200\Bar*
  \item Hauptventil mit Standardgewinde
  \item Daran Druckminderer/Berieselungseinheit
  angeschlossen
\end{itemize}
}
{}
{}
{}


\TextSlide{Berieselungseinheit}
{\index{Berieselungseinheit}
Die Berieselungseinheit besteht im Wesentlichen aus einem
\T{Druckminderer}. Der Druckminderer hat ein
\I{Abgabeventil}, welches den \E{Druck drosselt} und eine
Abgabe an den Patienten ermöglicht. Je nach Bauart kann der
Druckminderer über ein \E{Regelventil}, ein
\E{Flowmeter}\index{Flowmeter} und ein \E{Manometer}
verfügen.  Mit dem \T{Regelventil} kann man die Durchfluss-
bzw. Abgaberate ein\-stellen,  das \T{Flowmeter}
(Durchflussanzeige) zeigt die aktuelle Durchflussrate in 
Liter pro Minute (\LPerMin) an. Das \T{Manometer}  zeigt den
Druck in der Sauerstoffflasche an, wenn das Hauptventil
geöffnet ist.}
{
\begin{itemize}
  \item Druckminderer
  \item Regelventil
  \item Flowmeter \Sign{→} \LPerMin
  \item Manometer \Sign{→} \Bar
\end{itemize}
}
{}
{}
{}






\subsection{Umgang mit Sauerstoff und Druckflaschen}

\TextSlide{\TxGefahren}
{%
Sauerstoff ist \EE{brandfördernd} und zusammen mit Fett
sogar \Ca{selbstentzündlich}! Neben der Brand- und
Explosionsgefahr birgt der große Druck unter dem die Flaschen
stehen weitere Gefahren: Sollte ein Ventil abbrechen, kann die
Flasche eine enorme Kraft entwickeln und zu einer Rakete
werden. Die Kraft einer mit 200\Bar* gefüllten Flasche reicht
aus, um Mauern zu durchbrechen!


}{
\Ca{← Siehe Fließtext!}
}{}{}{}

\TextSlide{Beenden der \texorpdfstring{\ce{O2}}{O2}-Gabe und Wechsel von Druckflaschen}
{Nach der Verwendung des Sauerstoffsystems wird das
Hauptventil geschlossen und der \EE{Überdruck} aus den
Druckleitungen bzw. dem Druckminderer \EE{abgelassen}. Erst
danach wird das Abgabeventil ebenfalls verschlossen.

Gleiches gilt für den Wechsel einer Sauerstoffflasche: Es
muss -- zur eigenen Sicherheit -- sichergestellt sein, dass
\EE{im Entnahmeschenkel \Ca{kein Überdruck}} herrscht. Vor
dem Abmontieren des Druckminderers oder einer Druckleitung
muss dies auf alle Fälle \EE{extra am Manometer kontrolliert
werden}! Das \EE{Abgabeventil} muss beim Abmontieren
\EE{geöffnet} sein.



}{
\Ca{← Siehe Fließtext!}
}{}{}{}

\Text{Lagerung}
{
Sauerstoffflaschen werden in geeigneten
Haltevorrichtungen gelagert. Sie sollen nicht in der Sonne
oder im heißen Umfeld gelagert werden, da dadurch ein
falsch hoher Druck angezeigt wird.
}{}
{}{}{}

\Text{\Ca{Sicherheitshinweise}}
{\Cave{\Ca{Achtung! Vorsicht beim Hantieren -- Zerknall- und
Explosionsgefahr!}
\begin{enumerate}
    \item Sauerstoff ist zusammen mit Fett selbstentzündlich!
    \begin{enumerate}
        \item Die Armaturen und Gewinde sind fettfrei zu
        halten!
        \item Ventile dürfen nicht mit fetten oder öligen
        Händen bedient werden!
    \end{enumerate}
    \item Das Hantieren mit Feuer und offenem
    Licht ist in der Nähe von Sauerstoffflaschen verboten!
    \item Die Druckflaschen sind immer zu sichern und vor
    Umfallen, etc. zu schützen!
\end{enumerate}

\Ca{Lebensgefahr!} Vor der Abmontage eines
Druckminderers oder einer Druckleitung muss die
Überdruckfreiheit geprüft und das Abgabeventil geöffnet
werden!

}
\Fazit{Die Sauerstoffflaschen sind pfleglich zu behandeln!}}
{}
{}
{}
{}

\subsection{Heiteres Rechnen mit Sauerstoff}

Die insgesamt verfügbare Menge wird errechnet, indem man
die Füllgröße der Flasche (in Liter) mit dem
Flaschendruck (in Bar) multipliziert. Zur Berechnung, wie
lange die vorhandene Sauerstoffmenge bei einer bestimmten
Abgaberate reicht, \EE{dividiert man die gesamt verfügbare
Menge durch die Abgaberate}.
Zu Bedenken ist dabei, dass in der Flasche  mind.
\EE{10\Bar* Restdruck} verbleiben sollen und eine
\EE{ausreichende Zeitreserve} einzuplanen ist!  

{\FontSansNarrowFamily
\begin{equation}
\text{Verfügbares Volumen in \Liter} = {\begin{array}{c}
\text{Volumen der Flasche in \Liter}\\
\text{(Füllgröße)} \end{array}}
\times 
\begin{array}{c}{\text{Druck in bar}}\\
\text{(Flaschendruck)}
\end{array}
\end{equation}

\begin{equation}
\text{Zeit} = \frac{\text{Verfügbare
Menge}}{\text{Abgaberate}}
\end{equation}
}
% \begin{equation}
% V_\sum = V_F \times p
% \end{equation}
% 
% \begin{equation}
% t = \frac{V_\sum}{Flow}
% \end{equation}


%oder kurz: 

%\begin{equation}
%l = {Volumen} \times {Druck}
%\end{equation}





\Fazit{Es ist ein \EE{Restdruck} abzuziehen und eine
ausreichende \EE{Zeitreserve} einzuplanen!}



%oder kurz: 

%\begin{equation}
%min = \frac{l}{l / min}
%\end{equation}


\minisec{Beispiel}

\begin{quote}
Man hat eine Flasche mit einer Füllgröße von 10\Liter*. Diese
steht unter einem Druck von 150\Bar*: Daraus ergibt
sich, dass in der Flasche 1\,500\Liter* \ce{O2} sind.

Der Patient benötigt 5\,\unitfrac{\Liter}{\Min} \ce{O2} : 

Der Inhalt würde theoretisch für 300 Minuten reichen. Der
Restdruck sowie eine Zeitreserve sind dabei
aber noch \E{nicht} berücksichtigt!
\end{quote}

\minisec{Beispiel}

\begin{quote}
Gegeben ist: 
\begin{itemize}
  \item \ce{O2}-Flasche mit Füllgröße von 10\Liter*, Druck von
160\Bar*
\item Patient: \ce{O2}-Bedarf von 5\LPerMin*
\item Fahrtdauer laut Routenplaner:
4\Hour*~58\Min*
\end{itemize}

Nach Abzug des Restdruckes ergibt sich ein nutzbarer Druck
von 150\Bar* und folglich eine nutzbare \ce{O2}-Menge von
1\,500\Liter*. Der Patient könnte daher für 300\Min* mit
Sauerstoff versorgt werden. Die (geschätzte) Fahrtdauer
von 4\Hour*~58\Min* entspricht 298\Min*, es ergibt sich
eine Zeitreserve von 2\Min*.

Da die Zeitreserve von 2\Min* zu kurz ist, kann der
Transport mit dieser Ausstattung \E{nicht} durchgeführt
werden. 
\end{quote}

{\TakeNotesBoxLarge}

\subsection{Verabreichung von Sauerstoff}
\index{Sauerstoff!Verabreichung}


\Layout{60}{Einleitung}
{%
Die Verabreichung von Sauerstoff ist eine wichtige und häufige
Maßnahme. Sauerstoff kann sowohl einem selbst atmenden Patienten
zusätzlich gegeben werde (\T{Berieselung}), aber auch bei der
\EE{Beatmung} zugeführt werden. 

\bigskip

}
{}
{./images/motal-aass/o2_4765-00800.jpg}
{}
{Motal}


\Fazit{Unterscheide Berieselung (erhaltene Eigenatmung) und Beatmung!}

\subsubsection{Sauerstoff-Berieselung}
\index{Berieselung}
\index{Berieselungs!-maske}
\index{Berieselungs!-brille}
\label{topic:berieselung}
\label{topic:berieselungsbrille}
\label{topic:berieselungsmaske}

Die \ce{O2}-Berieselung kann \EE{nur bei vorhandener
Spontanatmung} durchgeführt werden! Zum Einsatz kommen
Berieselungsbrillen, Berieselungsmasken und
Berieselungsmasken mit Reservoir\ZFootnote{In Einzelfällen
kann auch mit einem Beatmungsbeutel berieselt werden, dadurch
sind Sauerstoffkonzentrationen bis zu 100\PC* möglich.
Allerdings ist die Anwendung komplizierter (Maske muss
ständig dicht angelegt werden) und die Materialkosten sind
deutlich höher (Desinfektion notwendig, bzw.
Neuanschaffung bei Einwegprodukten)}.

% \marginpar{\T{Sauerstoff\-brille}}
% \marginpar{\T{Sauerstoff\-maske}}
% \marginpar{\T{Sauerstoff\-maske + Reservoir}}

% zwei kurze, in die Nasenöffnungen  eingeführte Kunst\-stoff\-stutzen
% 
% bei ca. 4\,l O2 /min bis 40\% \ce{O2}  in Atemgas  möglich
% 
% 
% Bei behinderter Nasenatmung kann die  Sauerstoffbrille im Mundwinkel
% plaziert werden. 

\begin{PictureSeriesWide}{}
{Bilderserie: Mittel zur Berieselung mit Sauerstoff. \SourcePicture{Ch. Pallinger}}
\PictureSeriesRowWide{3}
{./images/pallinger-christoph-aass/Anton_32988_v2-AASS-0112mm.jpg}
{Sauerstoffbrille}
{./images/pallinger-christoph-aass/Anton_32987_v2-AASS-0112mm.jpg}
{Sauerstoffmaske}
{./images/pallinger-christoph-aass/Anton_32986_v2-AASS-0112mm.jpg}
{Sauerstoffmaske mit Reservoir}
{}
{}
\end{PictureSeriesWide}


\Layout{57}
{}
{}
{}
{\begin{tabulary}{\linewidth}[H]{LLLL}
\addlinespace\toprule\addlinespace
\noindent
 &
\noindent\TabHeading{Indikation}
 &
\noindent\TabHeading{Vorteile}
&
\noindent\TabHeading{Nachteile}
\\\addlinespace\midrule\addlinespace 
\TabSub{Sauerstoff\-brille}
 &
Spontanatmung

\ce{O2}-Flow \EE{bis~5\LPerMin*}
 &
gute Toleranz

fehlendes Engegefühl

Sprechen \& Husten möglich
 &
ungenaue Dosierung

Austrocknung  der Schleimhäute 

niedriger Flow

Nasenatmung Voraussetzung, ungeeignet bei Verlegung
\\\addlinespace\midrule\addlinespace
\TabSub{Sauerstoff\-maske}
 &
Spontanatmung

\ce{O2}-Flow \EE{ab~5\LPerMin*}
 &
höhere \ce{O2}-Konzentration

Atemkontrolle erleichtert (Beschlagen der Maske)
 &
\ce{CO2}-Rückatmung  wenn \ce{O2}-Durchfluss
{\textless}~5\LPerMin*!) 

Abflussbehinderung (Erbrechen, \ldots)

für Patienten evtl. unangenehm bzw. gewöhnungsbedürftig
\\\addlinespace\midrule\addlinespace
\TabSub{Sauerstoff\-maske mit Reservoir}
 &
Spontanatmung

\ce{O2}-Flow ab 5\LPerMin*
 &
Wie oben

Noch höhere \ce{O2}-Konzentration
 &
Wie oben

Reservoir muss zuerst händisch befüllt werden.
\\\addlinespace\midrule\addlinespace
\TabSub{Beatmungsbeutel mit Reservoir}
 &
Spontanatmung, assitierte und kontrollierte Beatmung
 &
Bis zu 100\PC \ce{O2}-Konzentration

Berieselung und Beatmung möglich
 &
deutlich aufwändiger und höhere Materialkosten

Einzelfälle!
\\\addlinespace\bottomrule
\end{tabulary}}
{Verabreichungsarten von Sauerstoff}
{}




\MassnahmenMK{Spezielle Maßnahmen}{Sauerstoffberieselung}
{m:sauerstoffberieselung}
{%
Bei jedem Patienten, bei dem eine
lebensbedrohliche Störung einer vitalen Funktion eingetreten
ist oder einzutreten droht (\enquote{Notfallpatient}), soll,
sofern keine Kontraindiktaionen vorliegen, soviel Sauerstoff
verabreicht werden, dass die Sauerstoffsättigung (Sp\ce{O2})
im Bereich von \EE{\VonBis{94}{98}\PC*} liegt.
% 
% \begin{enumerate}
%   \item Kontraindikationen und Gegenanzeigen prüfen:
%   \begin{itemize}
%     \item COPD (\FullPrettyRef{topic:copd})
%     \item Hyperventilationssyndrom, Hyperventilationstetanie
%     (\FullPrettyRef{topic:hyperventailationssyndrom})
%   \end{itemize}
%   \item Situationsgerechte Dosierung je nach
%   zugrundeliegender Erkrankung. Grundsätzlich soll ein
%   Sp\ce{O2} von
% \EE{\VonBis{94}{98}\PC*} erreicht werden.
%   \item Auswahl des Hilfsmittels:
%   \begin{itemize}
%     \item \E{Sauerstoff\-brille}
%     \ce{O2}-Flow bis~5\LPerMin*,
%     \item \E{Sauerstoff\-maske}
%     \ce{O2}-Flow ab 5\LPerMin*,
%     \item \E{Sauerstoff\-maske mit
%     Reservoir} \ce{O2}-Flow ab 5\LPerMin*,
%   \end{itemize}
%   \item Patientenaufklärung
%   \item Sauerstoffgerät einschalten und Durchflußrate
%   einstellen
%   \item Bei Verwendung einer Sauerstoff\-maske mit
%   Reservoir: Reservoir füllen
%   \item Sauerstoffbrille/-maske positionieren 
% \end{enumerate} 
% 
% \cite{AsboeLehrmeinungRd2011-09}
}


\subsubsection{Sauerstoffzufuhr bei der Beatmung}

\TextSlide{\TxBeschreibung}{%
Bei der Beatmung soll Sauerstoff zugeführt werden. Beim
Beatmungbeutel schließt man dazu einen
\ce{O2}-\EE{Verbindungsschlauch} an den Druckminderer und
den Beatmungsbeutel an. Wenn vorhanden soll ein
\ce{O2}-\EE{Reservoir} verwendet werden, dadurch wird die
Sauerstoffkonzentration auf fast 100\PC* erhöht.
Bei Beatmungsgeräten kann -- je nach Typ -- das
Mischungsverhältnis von Sauerstoff und Umgebungsluft
eingestellt werden.\ZFootnote{Z.\,B. Schalthebel
\enquote{Air Mix} bei Weinmann Medumat{\texttrademark} Standard.}
}{%
\begin{itemize}
  \item Beamtungsbeutel:
  \begin{itemize}
    \item \ce{O2}-\EE{Verbindungsschlauch}
    \item \ce{O2}-\EE{Reservoir}
  \end{itemize}
  \item Beatmungsgerät: einstellbar
\end{itemize}
}{}{}{}














% 
% \needspace{0.5\textheight}
% ~

\section{Beatmung}
\label{topic:beatmung-hilfsmittel}
\label{topic:beatmung-assistiert}
\label{topic:beatmung-kontrolliert}
\index{Beatmung}
\index{Beatmung!Grundsätzliches}

\TextSlide{\TxBeschreibung}
{Bei \EE{fehlender} oder \EE{nicht ausreichender
Spontanatmung} (Eigenatmung \textless\,8 bzw.
\textgreater\,35\PerMin* oder Atemzugvolumen nicht
ausreichend) wird eine \EE{künstliche Beatmung}
durchgeführt.
Dabei wird Eigenatmung des Patienten künstlich durch geeignete Hilfsmittel ersetzt oder ergänzt.
Es handelt es sich um eine
\EE{Überdruckbeatmung}\ZFootnote{Theoretisch kann auch
mittels einer Unterdruckkammer eine Unterdruckbeatmung
durchgeführt werden (\I[noindex]{Eiserne Lunge}), dies ist
aber sehraufwendig und wird praktisch nicht durchgeführt.},
d.\,h. es wird mit Druck Luft in die Atemwege des Patienten
gepumpt. (Im Unterschied dazu wird bei der natürlichen
Atmung Luft mit Unterdruck in den Patienten eingesogen.) Man
unterscheidet zwischen einer \E{kontrollierten} und einer
\E{assistierten} Beatmung.

Bei der \T[Beatmung!kontrollierte]{kontrollierten Beatmung}
wird die gesamte Atmung künstlich aufrecht erhalten;
hingegen bei der \T[Beatmung!assistierte]{assistierten
Beatmung} hat der Patient eine eigene Spontanatmung, die
aber nicht ausreichend ist (zu langsam oder zu seicht). Man
unterstützt dabei die Atembemühungen des Patienten, indem
man zusätzliche oder tiefere Atemhübe verbreicht.

Eine Beatmung durch Fachpersonal sollte wenn möglich immer
mit Hilfsmitteln wie dem \E{Beatmungsbeutel} oder einem
\E{Beatmungsgerät} durchgeführt werden.
}{
\begin{itemize}
  \item Künstliche Überdruckbeatmung
  \item Kontrollierte Beatmung: Atmung wird vollständig
  übernommen
  \item Assistierte Beatmung: Eine vorhandene, aber
  unzureichende Spontanatmung wird unterstützt
\end{itemize}
}{}{}{}

\TextSlide{\TxGefahren}
{Der notwendige \EE{Überdruck} kann zu Problemen führen, wie
z.\,B. Überschreiten des \E{Speiseröhren-Öffnungsdrucks}
(\EE{Magenbeatmung}, \EE{Magenblähung}, \EE{Erbrechen}) oder
Lungenschäden.
Es ist oft zu beobachten, das (wahrscheinlich in Folge von
Stress) Beatmungen zu schnell bzw. zu tief erfolgen. Dies
kann katastrophale Auswirkungen haben, da es zu einer
\EE{Hyperventilation} und damit zu einer atmungsbedingen
Alkalose kommt (Störung des Säure-Basen-Haushaltes, vgl.
\prettyref{topic:alkalose-respiratorische}).

\Cave{Cave: Hyperventilation und atmungsbedingen Alkalose}

\Fazit{Faustregel: Eine kontrollierte Beatmung eines
Erwachsenen soll grundsätzlich im Eigenrhythmus des Helfers
erfolgen: Helfer atmet \Sign{→} Patient wird beatmet.}
}{
\begin{itemize}
  \item Zu hoher Druck \Sign{→} Magenbeatmung/-blähung,
  Erbrechen
  \item Zu schnell, zu tief: Hyperventilation,
  atmungsbedingte Alkalose
\end{itemize}
}{}{}{}

\TextSlide{Hilfsmittel}
{Eine Beatmung wird normalerweise mit Hilfsmitteln
durchgeführt. Helferseitig kommen dabei ein 
\EE{Beatmungsbeutel} (\prettyref{topic:beatmungsbeutel})
oder ein \EE{Beatmungsgerät} zum Einsatz. 
Patientenseitig werden \EE{Beatmungsmasken} oder
\EE{Beatmungsschläuche} (Endotrachealtuben,
Larynxtuben, \Z{Larynxmasken, Combituben, \ldots})
angewendet. Um die Atemwege freizuhalten, kann
u.\,U. unterstützend ein \T{Guedel-Tubus} angwendet werden.

\bigskip

\bigskip
}{
\begin{itemize}
  \item Helferseitig:
  \begin{itemize}
    \item Beatmungsbeutel
    \item Beatmungsmaske
  \end{itemize}
  \item Patientenseitig
  \begin{itemize}
    \item Beatmungsmaske
    \item Beatmungsschlauch
  \end{itemize}
  \item Freihalten der Atemwege: Guedel-Tubus
\end{itemize}
}{}{}{}

\begin{table}[H]
\Captionof{table}{Mögliche Sauerstoffkonzentrationen bei der Beatmung.}
\FontTableBase
\begin{tabulary}{\linewidth}{Lc}\addlinespace\toprule\addlinespace
%\addlinespace 
\noindent\TabHeading{Art der Beatmung} &
\noindent\TabHeading{\ce{O2}-Konzentration}
\\\addlinespace\midrule
Mund-zu-Mund Beatmung
&
\Range{16}{17}\PC*
\\\addlinespace
Beutel-Masken-Beatmung
&
21\PC*
\\\addlinespace
Beutel-Masken-Beatmung mit angeschlossenem \ce{O2} (15\LPerMin*)
&
\VonBis{40}{70}\PC*
\\\addlinespace
Beutel-Masken-Beatmung mit angeschlossenem \ce{O2} (15\LPerMin*) und
Reservoir &
fast 100\PC*
\\\addlinespace\bottomrule
\end{tabulary}

\end{table}



\subsection{Einfache Atemwegssicherung}

Sobald eine Atemwegsverlegung bemerkt wird, müssen sofort
Maßnahmen eingeleitet werden, die die Atemwege öffnen und
offen halten. Es gibt 3 Maßnahmen, die die Atemwege im Fall
einer Atemwegsverlegung -- verursacht durch die Zunge oder
davon oberhalb befindlicher Strukturen -- frei halten
können \cite{Erc2010Sec4}:

\begin{itemize}
  \item \EE{Überstrecken} des Kopfes 
  \item Anheben des \EE{Kinns}
  \item Vorschieben des Unterkiefers
  mittels \EE{Esmarch-Handgriff}
\end{itemize}

\Text
{Kopf Überstrecken, Kinn anheben}
{Zum \IX[Ueberstrecken@Überstrecken]{Überstrecken} des Kopfes und Anheben wird die Hand
auf die Stirn des Patienten gelegt und der Kopf
leicht überstreckt, dabei werden die Fingerspitzen
der anderen Hand unterhalb des Kinns platziert und heben
dieses leicht an. \cite{Erc2010Sec4}}
{}
{}
{}
{}



\TextSlide
{Esmarch-Handgriff}
{Der \IX{Esmarch-Handgriff} verstärkt den Effekt des
Überstreckens des Kopfes. Der Helfer befindet sich hinter
dem Kopf des Patienten und platziert seine Finger
\VonBis{2}{5} am Ende des Unterkiefers am Kieferwinkel und
den Daumen am vorderen Bereich des Unterkiefers. Durch \E{auf-}
und \E{vorwärtsgerichteten} Druck der Finger am
Kieferwinkel kann der Unterkiefer nach oben und vorn
geschoben werden, der Daumen öffnet durch \E{sanften Druck
auf das Kinn} vorsichtig den Mund. \cite{Erc2010Sec4}

}
{
\begin{itemize}
  \item Finger
\VonBis{2}{5} am Kieferwinkel

$\rightarrow$ aufwärts, vorne
\item Daumen seitlich am Kinn

$\rightarrow$  sanfter Druck nach unten
\end{itemize}
}
{}
{}
{}

\Text
{Guedel-Tubus}
{\index{Guedel-Tubus}
Der Guedeltubus ist ein bogenförmiges, abgeflachtes
Plastikrohr, welches zum Offenhalten der Atemwege bei tief
bewusstlosen Patienten eingesetzt werden kann. Er wird über
den Mund eingeführt und kommt mit seinem Ende im Rachen zu
liegen. Er erleichtert die Beatmung, da er die Zunge
aus dem Weg räumt, bietet jedoch \E{keinen}
Aspirationschutz. \FullPrettyRef{fig:guedel-einfuehren}
zeigt das Einführen und die Lage des Guedel-Tubus.
}
{}
{}
{}
{}

\Customized{\begin{description}
  \item[{\CustAsboe}] Die Anwendung eines
  Güdel-Tubus ist für die Beatmung (nur) während der
  Reanimation ab der Ausbildungsstufe Rettungssanitäter freigegeben.
\cite{AsboeLehrmeinungRd2011-15}
\end{description}}



\begin{PictureSeriesWide}{}{Einführen des Guedel-Tubus
\SourcePicture{Hirtler}\label{fig:guedel-einfuehren}}
\PictureSeriesRowWide{3}
{./images/hirtler-aass/GuedelTubus1}
{Abmessen, \ldots}
{./images/hirtler-aass/GuedelTubus2}
{\ldots\ zuerst wird der Guedel-Tubus \enquote*{verkehrt}
herum eingeführt \ldots}
{./images/hirtler-aass/GuedelTubus3}
{\ldots\ und anschließend mittels Drehbewegung endgültig
positioniert.}
{empty}
{}
\end{PictureSeriesWide}

\begin{PictureSeriesWide}{}{Übersicht: Lage verschiedener
Hilfsmittel zur Atemwegssicherung
\SourcePicture{Hirtler}\label{fig:uebersicht-lage-atemwegshilfen}}
\PictureSeriesRowWide{3}
{./images/hirtler-aass/PositionGuedel-color}
{Position des Guedel-Tubus im Querschnitt}
{./images/hirtler-aass/LageLarynxTubus}
{Im Vergleich: Lage eines Larynx-Tubus \ldots} 
{./images/hirtler-aass/LageTubus-color}
{\ldots\ und eines endotrachealen Tubus.}
{empty}
{}
\end{PictureSeriesWide}

{\TakeNotesBoxLarge}

\subsection{Beatmungsbeutel}
\label{topic:beatmungsbeutel}
\index{Beatmungs!-beutel}



\Text{\TxBeschreibung}
{Der selbstfüllende Beatmungsbeutel ist das einfachste
Hilfsmittel, welches dem Fachpersonal zur Verfügung steht, und kommt in
vielen Bereichen (Rettungsdienst, Krankenhaus, \ldots) zum
Einsatz. Er ermöglichst die kontrollierte oder
unterstützende Beatmung eines Patienten. Er kann mittels
einer Gesichtsmaske (Beatmungsmaske) angewendet, oder an
andere Systeme (Tuben) angeschlossen werden. }
{%
\begin{itemize}
  \item Hilfsmittel zur kontrollierten und unterstützenden Beatmung
\end{itemize}
}
{}
{}
{}

\Layout{37}{Abbildung}
{}
{}
{./images/pallinger-christoph-aass/Ambubeutel_32741-AASS-0176mm.jpg}
{Selbstfüllender Beatmungsbeutel
\enquote{Ambu{\texttrademark} Mark IV} mit Reservoir, \ce{O2}-Verbindungsschlauch,
Bakterienfilter, Beatmungsmaske und aufgestecktem PEEP-Ventil.} {Ch.
Pallinger}

% \Layout{22}{Bestandteile}
% {
% Ein Beatmungsbeutel besteht im Idealfall aus folgenden Teilen:
% \begin{itemize}
% 	\item Beatmungsmaske 
% 	\item Bakterienfilter
% 	\item Ein-/Ausatemventil mit angeschlossenem PEEP-Ventil
% 	\item Beatmungsbeutel
% 	\item Reservoir, welches mit einem Verbindungsschlauch
% 	(\E{\ce{O2}-Line}) mit einer Berieselungseinheit verbunden
% 	ist.
% \end{itemize}
% }{}{}{}

\TextSlide{Bestandteile}{%
Ein Beatmungsbeutel besteht üblicherweise aus einem
\EE{Luftzufuhrventil}, einem \EE{Balg} und einem
\EE{T-Stück} mit einem Ventil.
An der Luftzufuhr kann eine \EE{Sauerstoffleitung} und
ein \EE{Sauerstoffreservoir} angeschlossen werden, um den
Sauerstoffgehalt des Einatemgases zu erhöhen. Am einfachsten
erfolgt dies über einen Verbindungsschlauch und eine
\ce{O2}-Berieselungseinheit, welche mit einem Flow von 15\LPerMin*{}
eingestellt ist. Mit Hilfe eines Reservoirs, welches
Sauerstoff während der Beatmungspausen zwischenspeichert,
kann die \ce{O2}-Konzentration des Einatemgases weiter
gesteigert werden.~\ZFootnote{Andere \ce{O2}-Systeme sind am
Markt verfügbar (Demand-Ventile, etc.), aber nur begrenzt
verbreitet.}

\Fazit{Wenn vorhanden, soll der Beutel mit einer
Sauerstoffzufuhr verbunden werden.}

Am T-Stück gibt es \E{3 Schenkel}: Einer kommt vom
Beatmungsbalg, einer führt zum Patienten und einer zum
Auslass. Das \EE{Ventil} des T-Stückes \E{trennt die Ein- von
der Ausatemluft} des Patienten. Am Patientenschenkel soll ein
\EE{Bakterienfilter} und Messvorrichtungen (z.\,B. zur
\ce{CO2}-Messung) angebracht werden. Er wird mit einer
\EE{Beatmungsmaske} oder einem Tubus -- je nach Bedarf --
verbunden. Am Ausführungsschenkel kann ein \EE{PEEP}-Ventil
aufgesteckt werden.
}{%
\begin{itemize}
  \item Luftzufuhrventil
  \begin{itemize}
    \item \ce{O2}-Leitung
    \item Reservoir
  \end{itemize}
  \item Balg
  \item T-Stück
  \begin{enumerate}
    \item Vom Balg kommend
    \item Zum/vom Patienten: + Bakterienfilter,
    Beatmungsmaske od. Tubus
    \item Ausführung, evtl. mit PEEP-Ventil
  \end{enumerate}
\end{itemize}
}
{./images/geraete/ambu-komplett-zerlegt.jpg}
{Bestandteile eines
Beatmungsbeutels.}
{}




% \Layout{57}{}{}{}{
% \begin{tabularx}{\linewidth}{XXX}
% {\includegraphics[width=\linewidth]{./images/pallinger-christoph-aass/Ambubeutel_32741-AASS-0112mm.jpg}
% \Captionof{figure}{Beamtungsbeutel \enquote{Ambu{\texttrademark} Mark
% IV} mit Reservoir, \ce{O2}-Verbindungsschlauch,
% Bakterienfilter und aufgestecktem PEEP-Ventil.
% \SourcePicture{Ch. Pallinger}}} & {\includegraphics[width=\linewidth]{}
% \Captionof{figure}{}} &
% {\includegraphics[width=\linewidth]{}
% \Captionof{figure}{}}
% \\\addlinespace\bottomrule
% \end{tabularx}
% }
% {.}
% {}

\begin{PictureSeries}{}{Bilderserie: Beatmungsmasken}
\PictureSeriesRow{2}
{./images/pallinger-christoph-aass/Beatmungsmaske_32906_v2-AASS-0112mm.jpg}
{Eine \E{Beatmungs}maske. \SourcePicture{Ch. Pallinger}}
{./images/pallinger-christoph-aass/Mundstuecke_32723-AASS-0112mm.jpg}
{Beatmungsmasken unterschiedlicher Größen:
0, 2, 3/4, 5. Die Farben der Masken haben
keine besondere Bedeutung. \SourcePicture{Ch. Pallinger}}
{empty}
{}
{empty}
{}
\end{PictureSeries}


\TextSlide{Hygienische Wiederaufbereitung}
{%
Es gibt Beatmungsbeutel als Einweg- und Mehrwegprodukt.
Einwegprodukte müssen nach Gebrauch entsorgt werden.
Mehrwegprodukte können durch Desinfektion gem.
des jeweils gültigen Hygieneplans wiederaufbereitet werden.
I.\,d.\,R. muss der
Beatmungsbeutel dazu \EE{komplett zerlegt} werden!
Die Herstellerhinweise sind zu beachten.
}{
\begin{itemize}
  \item Ein- und Mehrwegprodukte
  \item Ggfs. Desinfektion gem. Hygieneplan
\end{itemize}


}{}{}




% \subsubsection{Technik der Beutel-Masken-Beatmung}

\TextSlide
{Technik der Beutel-Masken-Beatmung}
{Bei der Beutel-Masken-Beatmung erfolgt die Beatmung mittels
eines Beatmungsbeutel und einer Beatmungsmaske in geeigneter
Größe, welche dicht sitzend angelegt wird.
Der Kopf des Patienten wird
\EE{überstreckt}\footnote{Überstrecken des Kopfes bei der
Beatmung: Wenn der Atemweg durch andere Maßnahmen
(Guedel-Tubus, Wendel-Tubus, etc.; ärztliche Maßnahme!)
freigehalten wird, entfällt das Überstrecken des Kopfes.)}
und die Maske wird mittels \T{C-Griff} angelegt: Der Daumen
und der Zeigefinger umschließen die Beatmungsmaske und
bilden dabei die Form des Buchstaben \enquote*{C}, die
restlichen drei Finger heben das Kinn am Unterkiefer an.

Ein erwachsener Patient wird mit mit einem Volumen von ca.
\EE{500\Ml*} beatmet, welches über 1\Sec* verabreicht werden
soll (\E{nicht} stoßweise!). Die Beatmung soll grundsätzlich
in einer normalen Atemfrequenz
(\enquote*{\EE{Eigenfrequenz}}) erfolgen.
% Wird zu schnell (und
% damit zuviel) beatmet kann es zu schweren Störungen im Säure-Basen-Haushalt und in Folge zu
% einer schlechten Versorgung mit Sauerstoff im Körper kommen!
}
{\begin{itemize}
  \item Kopf überstrecken
  \item C-Griff
  \item ca. \EE{500\Ml*} über 1\Sec*
  \item Eigenfrequenz!
\end{itemize}

}
{}
{}
{}




\begin{PictureSeriesWide}{}{Die Technik der Beutel-Masken-Beatmung \SourcePicture{Hirtler}}
\PictureSeriesRowWide{3}
{./images/hirtler-aass/CGriff}
{C-Griff \ldots}
{./images/hirtler-aass/Beutelbeatmung1}
{\ldots\ und überstreckter Kopf}
{./images/hirtler-aass/Beutelbeatmung2}
{Normalerweise wird bei der Überdruckbeatmung mit einem
Beatmungsbeutel Luft in die Lunge gepumpt. Übersteigt der
Druck den Öffnungsdruck der Speiseröhre (ca.
20\,{\FontUnits mbar}), wird Luft auch in den Magen gepumpt.
}
{}
{}
\end{PictureSeriesWide}


% \subsubsection{Besonderheit: Beatmung während der
% Reanimation}

\Text
{Besonderheit: Beatmung während der
Reanimation}
{Die Beatmung während der Reanimation ist eine kontrollierte
Beatmung. Wird eine Beatmungsmaske verwendet, so ist das
Verhältnis zur Herzdruckmassage zu beachten (vgl.
\CO{[BLS, EHI]}; \prettyref{chp:BLS}, \ref{chp:EHI}). Sobald
die Beatmung über einen Beamtungsschlauch (Tubus) erfolgt,
kann sie unabhängig von der Herzmassage erfolgen.}
{}
{}
{}
{}










\subsection{Beatmungsgeräte}
\index{Beatmungs!-gerät}
\index{Beatmungs!-gerät!Medumat{\texttrademark}}
\index{Medumat{\texttrademark}}
\index{Medumat{\texttrademark}!Standard a}
\index{Medumat{\texttrademark}!Compact}
\index{Medumat{\texttrademark}!Standard}
\index{Medumat{\texttrademark}!Transport}
\label{topic:beatmungsgeraete}

\begin{figure}[H]
\begin{center}
\includegraphics[width=7.189cm,height=4.94cm]{./images/beatmungsgeraet-medumat.png}
\Captionof{figure}{Beatmungsgerät \emph{Medumat
Standard\texttrademark}. \SourcePicture{ASB
Floridsdorf-Donaustadt}}
\end{center}
\end{figure}




Mit dem
Beatmungsgerät \emph{{Medumat{\texttrademark}}} der
Fa.
Weinmann kann ein Patient sowohl beatmet als auch berieselt werden
(sofern ein Berieselungsmodul vorhanden ist). 
Für die Berieselung und Beatmung gibt es je ein eigenes
Bedienfeld. Die Beatmung mittels Beatmungsgerät ist
grundsätzlich dem ärztlichen Personal vorbehalten, die
Bedienung des Beatmungsbedienfeldes darf für nichtärztliches
Sanitätspersonal nur nach Anweisung eine Arztes
erfolgen.\ZFootnote{\EE{Bedienung des Beatmungsrätes durch
nichtärztliches Personal}: Ausdrückliche Ausnahmen können
für höher qualifiziertes Personal bestehen.}


{\TakeNotesBoxLarge}
% % \opt{ORGASBFLO}{
% % \begin{minipage}{\linewidth}
% \Customized{
% \begin{description}
% \item[\CustAsbFlo] Es kommen folgende Geräte
% zum Einsatz:
% \begin{itemize}
%   \item \EE{Notfallbeatmungsgeräte}:
%   \begin{itemize}
%     \item Hersteller: Fa. Weinmann, Produktfamilie:
%     \T{Medumat}:
%     \begin{inparaitem}
%       \item Medumat{\texttrademark} Compact
%       \item Medumat{\texttrademark} Standard
%       \item Medumat{\texttrademark} Standard a
%     \end{inparaitem}
%   \end{itemize}
%   \item \EE{(Intensiv-)Transportbe\-atmungs\-ge\-rät}\nocite{ThierbachInterhospitaltransfer1}:
%   \begin{itemize}
%     \item Hersteller: Fa.
%     Weinmann, Produktfamilie: \T{Medumat{\texttrademark}}:
%     \begin{inparaitem}
%       \item Medumat{\texttrademark} Transport
%     \end{inparaitem}
%   \end{itemize}
% \end{itemize}
% \end{description}
% }

% 
% \end{minipage}
% }



\begin{figure}[H]
\includegraphics[width=\linewidth]{./images/geraete/ger-medumat-standard-bedienfeld.jpg}

\CaptionofDiff{figure}{Bedienfelder Berieselungseinheit \E{Modul Oxygen} (li.)
und Beatmungsgerät \E{Medumat{\texttrademark} Standard} (re.).

\E{Modul Oxygen (li.):} Flowmeter, Ein-Ausschalter und
Regelventil. Der Anschluss links dient der Verbindung der
Einheit mit dem Sauerstoffnetz des Fahrzeuges.

\E{Medumat{\texttrademark} Standard (re.):}
\Z{Air-Mix-Schalter, Ein-Ausschalter, Stellräder für
Atemfrequenz, Atemminutenvolumen (!) und Druckbegrenzung,
Beatmungsdruckanzeige, Warnleuchten für Verlegung
(\enquote*{Stenosis}), Lösung des Beatmungsschlauches
(\enquote*{Disconnection}), niederen Flaschendruck (\enquote*{<\,2.7\Bar*})
und schwache Batterie.} } {Bedienfelder Berieselungseinheit
\E{Modul Oxygen} (li.) und Notfallbeatmungsgerät
\E{Medumat{\texttrademark} Standard} (re.).}
\end{figure}



\subsection{PEEP und PEEP-Ventil}
% \index{PEEP-Ventil}
\label{topic:peep-ventil}

\Definition{\T{PEEP} steht für \enquote*{Positive End-Exspiratory
Pressure} (positiver end-exspiratorischer Druck), d.\,h.
positiver Druck in den Atemwegen am Ende der
Ausatmungsphase.}

\Layout{20}{\TxBeschreibung{}}
{%
Der \T{PEEP} sorgt dafür, dass am Ende der Ausatmungsphase
(Exspiration) der Druck in den Atemwegen nicht auf 0
absinkt, sondern stattdessen ein positiver Druck aufrecht
erhalten wird. Dadurch wird das Zusammenfallen der
Lungenbläschen (Alveolen, \prettyref{topic:alveolen})
verhindert.
Ein \Abk{PEEP} kann mittels eines
\T[PEEP-Ventil]{PEEP-Ventils} erreicht werden. Es wird dabei
auf die Ausatemöffnung des Beatmungsbeutels bzw.
Beatmungsschlauches aufgesteckt. Durch Drehen der Kappe kann
der gewünschte \Abk{PEEP} eingestellt werden. Die
Einstellung kann an der angebrachten \E{Skala} abgelesen
werden\ZFootnote{Der Druck wird je nach Modell in
\MetricUnit{cm} Wassersäule (\MetricUnit{cm\,\ce{H2O}}) oder
Millibar (\MetricUnit{mbar}) angebeben. Typische Werte sind
z.\,B. 5 oder 10\,\MetricUnit{cm\,\ce{H2O}}.}.
Es gibt \Abk{PEEP}-Ventil-Produkte sowohl als \EE{Einweg-},
als auch als \EE{Mehrweg}material!\ZFootnote{Moderne
Beatmungsgeräte verfügen oft über eine integrierte
PEEP-Funktion und benötigen kein zusätzliches
\Abk{PEEP}-Ventil.}
\Z{Bei der Maskenbeatmung ist zu beachten, dass der PEEP nur
erreicht werden kann, wenn das Beatmungssystem \E{luftdicht}
ist. Sobald die Maske nicht mehr dicht aufsitzt, verliert
das \Abk{PEEP}-Ventil seine Wirkung, da die Ausatemluft
unkontrolliert über das \E{Leck} entweichen kann.

}
}
{%
\begin{itemize}
  \item Positiver end-exspiratorischer Druck
  \item Zusammenfallen der
  Lungenbläschen verhindert
  \item Ärztliche Maßnahme
\end{itemize}
}
{./images/pallinger-christoph-aass/Peep_32909-AASS-0112mm.jpg}
{PEEP-Ventil}
{Ch. Pallinger}

\Cave{Die Anwendung des \Abk{PEEP} ist grundsätzlich eine
\EE{ärztliche Maßnahme}! }


\ChapterEnd


